\documentclass[11pt]{article}
\usepackage[a4paper,margin=27mm]{geometry}
\usepackage{amsmath,amssymb,amsthm,mathtools}
\usepackage[hidelinks]{hyperref}
\newtheorem{theorem}{Theorem}
\newtheorem{lemma}[theorem]{Lemma}
\title{Strict convexity of Gamma-ratio level curves}
\author{Research manuscript}
\date{October 28, 2026}
\begin{document}
\maketitle
\begin{abstract}
For every $c>1$, the equation
$\Gamma(h)\Gamma(h+2x)/\Gamma(h+x)^2=c$ defines a unique positive height
$h=h_c(x)$ for each $x>0$. We prove that $h_c$ is strictly convex on the
entire positive half-line. At $c=2$, this answers the question in Remark~16(a)
of Garrappa--Gerhold--Popolizio--Simon. The proof applies classical
perspective convexity to Euler's product for the Gamma function; strictness
follows from two consecutive factors of that product.
\end{abstract}

\section{The level curve}
\begin{theorem}\label{thm:main}
Let $c>1$. For each $x>0$ there is a unique $h_c(x)>0$ satisfying
\begin{equation}\label{eq:level}
 \frac{\Gamma(h_c(x))\Gamma(h_c(x)+2x)}{\Gamma(h_c(x)+x)^2}=c.
\end{equation}
For distinct $x_0,x_1>0$ and $0<\lambda<1$,
\[
 h_c\bigl(\lambda x_0+(1-\lambda)x_1\bigr)
 <\lambda h_c(x_0)+(1-\lambda)h_c(x_1).
\]
\end{theorem}

For $c=2$, the function in \eqref{eq:level} is the function $h$ of
\cite[Lemma~15]{GGPS}. That paper proves its strict convexity on
$[1,\infty)$ and proposes strict convexity on $(0,\infty)$ in
Remark~16(a). The remark is unchanged in the published version of record.
Its boundary classification for the Mittag--Leffler function in
Theorem~23 uses this same curve, including at parameters below one.
Theorem~\ref{thm:main} concerns the shape of this curve; it makes no claim
to resolve that paper's separate complete-Bernstein-function conjectures.

For $0<x<s$, put
\begin{equation}\label{eq:defect}
 F(x,s)=\log\Gamma(s-x)+\log\Gamma(s+x)-2\log\Gamma(s),
 \qquad g(u)=-\log(1-u^{-2})\quad(u>1).
\end{equation}
Euler's product for the Gamma function gives
\begin{equation}\label{eq:series}
 F(x,s)=\sum_{m=0}^{\infty}g\!\left(\frac{s+m}{x}\right).
\end{equation}
For clarity, the corresponding finite approximants for $\log\Gamma(z)$
are
\[
 z\log N+\log(N!)-\sum_{m=0}^{N}\log(z+m).
\]
In the combination \eqref{eq:defect}, both $\log N$ and $\log(N!)$
cancel. Each remaining summand is
$2\log(s+m)-\log(s-x+m)-\log(s+x+m)$, exactly the summand in
\eqref{eq:series}. Passing to the limit proves the identity.

All summands in \eqref{eq:series} are positive. For $m\geq1$,
\begin{equation}\label{eq:bound}
 0\leq g\!\left(\frac{s+m}{x}\right)
 \leq \frac{x^2}{(s+m)^2-x^2}\leq\frac{x^2}{m^2}.
\end{equation}
The first upper bound follows from $-\log(1-t)\leq t/(1-t)$ for
$0\leq t<1$. Thus the tail converges uniformly for fixed $x$ and $s>x$.
For fixed $x>0$, $F(x,s)$ is continuous and strictly decreasing in $s$.
Its limit is $+\infty$ as $s\downarrow x$, by the $m=0$ summand, and
zero as $s\to\infty$, by \eqref{eq:bound} and dominated convergence for
the series. Consequently there is a unique $s>x$ with $F(x,s)=\log c$.
Setting $h_c(x)=s-x$ proves existence and uniqueness in
Theorem~\ref{thm:main}.

\section{Perspective convexity and strictness}
Direct differentiation gives
\[
 g'(u)=-\frac{2}{u(u^2-1)},\qquad
 g''(u)=\frac{2(3u^2-1)}{u^2(u^2-1)^2}>0\quad(u>1).
\]
Hence $g$ is strictly convex. In particular,
$xg((s+m)/x)$ is a classical perspective of $g$ composed with an affine
shift. The perspective-convexity principle is standard
\cite[Propositions~2.3 and~2.13]{Combettes}; no new convexity principle is
needed here. The positive normalization by $x$ is relevant: $xF(x,s)$,
rather than $F(x,s)$ itself, is a sum of these convex perspectives.

We give the strict inequality explicitly. Let $x_0\ne x_1$ be positive,
$y_i=h_c(x_i)$, $s_i=x_i+y_i$, and write
\[
 a=\lambda,\quad b=1-\lambda,\qquad
 X=ax_0+bx_1,\quad S=as_0+bs_1,\qquad
 A=\frac{ax_0}{X},\quad B=\frac{bx_1}{X}.
\]
Here $A,B>0$, $A+B=1$, and $X<S$. For every $m\geq0$,
\begin{equation}\label{eq:weights}
 \frac{S+m}{X}
 =A\frac{s_0+m}{x_0}+B\frac{s_1+m}{x_1}.
\end{equation}
Convexity of $g$ therefore gives
\begin{equation}\label{eq:jensen}
 g\!\left(\frac{S+m}{X}\right)
 \leq A g\!\left(\frac{s_0+m}{x_0}\right)
      +B g\!\left(\frac{s_1+m}{x_1}\right).
\end{equation}
At least one of $m=0,1$ gives strict inequality. Indeed, equality for both
would force equality of both pairs of arguments, and subtraction yields
$1/x_0=1/x_1$, contrary to $x_0\ne x_1$.
The convergent sums in \eqref{eq:series} and the nonnegative termwise
gaps in \eqref{eq:jensen} imply
\[
 F(X,S)<A F(x_0,s_0)+B F(x_1,s_1)=\log c.
\]
On the other hand, $F(X,X+h_c(X))=\log c$. Strict decrease of $F(X,s)$
in its second variable gives $X+h_c(X)<S$. Subtracting $X$ gives
\[
 h_c(X)<a h_c(x_0)+b h_c(x_1),
\]
which completes the proof of Theorem~\ref{thm:main}.

\section{Formal and research status}
Independent source, mathematical, formal, manuscript and bounded novelty
reviews are complete, including reassessment of the explicit
classical-perspective reduction. The result is one family for all $c>1$;
the Euler product, perspective theorem, Jensen inequality and two-factor
strictness calculation are supporting tools, not separate outcomes.

The complete Lean certificate proves the literal Gamma-ratio statement
without residual analytic premises. In namespace \texttt{GammaLevel},
the endpoint is \texttt{source\_convexity}.
The level-$2$ specialization is \texttt{source\_remark16a}; existence and
uniqueness are separately certified by \texttt{existsUnique\_sourceLevel}.
A fresh isolated source replay passed for all nine modules and 52 public
declarations under Lean 4.34.1, Mathlib revision
\texttt{d13f23b723b}, and warnings as errors.
All audited declarations use only the standard axioms \texttt{propext},
\texttt{Classical.choice}, and \texttt{Quot.sound}.
A separate independent final review reproduced the extracted replay and
inspected every declaration type and axiom report. Numerical checks are
supplementary evidence only; the campaign acceptance record is separate.

The bounded literature search found no earlier resolution of this level-curve
question; worldwide priority is not asserted. The proof applies the classical
perspective method to the Gamma-ratio equation.

\begin{thebibliography}{9}
\bibitem{GGPS} R. Garrappa, S. Gerhold, M. Popolizio and T. Simon,
\emph{On some inequalities for the two-parameter Mittag--Leffler function
in the complex plane}, J. Math. Anal. Appl. \textbf{551} (2025), 129588.
\href{https://doi.org/10.1016/j.jmaa.2025.129588}{doi:10.1016/j.jmaa.2025.129588}.
See also arXiv:2410.11852v4, Lemma~15 and Remark~16(a).
\bibitem{Combettes} P. L. Combettes,
\emph{Perspective functions: properties, constructions, and examples},
arXiv:1610.01552v3 (2017), Propositions~2.3 and~2.13.
\end{thebibliography}
\end{document}
